\Article{348}
\textsc{Problema.} \emph{Data aequatione}
\[
x^m+Ax^{m-1}+Bx^{m-2}+Cx^{m-3}+\text{etc.}+M=0
\]
\emph{cuius radices sunt $x=a$, $x=b$, $x=c$ etc., invenire aequationem, cuius radices sint $x=a^\tau$, $x=b^\tau$, $x=c^\tau$ etc.}

\SourceLeaf{233}
\emph{Solutio prima.} Quaerantur per theorema notum summae radicum aequationis propositae, earum quadratorum, cuborum etc. usque ad potestatem $m\tau^{\mathrm{tam}}$. Hinc igitur habentur etiam summae radicum aequationis quaesitae nec non quadratorum etc. scilicet $\Sigma a^\tau$, $\Sigma a^{2\tau}$ etc., unde per idem theorema coëfficientes determinari possunt.

Ad praxin quidem haec solutio est facilior; sed ad institutum nostrum nec non ad ostendendum, coëfficientes aequationis quaesitae fore integros, si aequationis propositae coëfficientes fuerint integri, quae sequitur magis est accomodata.

\emph{Solutio secunda.} Sit $\theta$ radix prima aequationis $x^\tau=1$, fiatque productum ex
\[
\begin{aligned}
&x^m+Ax^{m-1}+Bx^{m-2}+\text{etc.}\\
&x^m+A\theta x^{m-1}+B\theta\theta x^{m-2}+\text{etc.}\\
&x^m+A\theta\theta x^{m-1}+B\theta^4x^{m-2}+\text{etc.}\\
&\text{etc.}\\
&x^m+A\theta^{\tau-1}x^{m-1}+B\theta^{2\tau-2}x^{m-2}+\text{etc.}
\end{aligned}
\]
Huius itaque producti radices erunt
\[
\begin{array}{llll}
a,&\theta a,&\theta\theta a&\text{etc.}\\
b,&\theta b,&\theta\theta b&\text{etc.}\\
c,&\theta c,&\theta\theta c&\text{etc.}
\end{array}
\]
i. e. productum aequale erit huic
\[
(x^\tau-a^\tau)(x^\tau-b^\tau)(x^\tau-c^\tau)\ldots
\]
adeoque huius formae
\[
x^{\tau m}+A'x^{\tau(m-1)}+B'x^{\tau(m-2)}+\text{etc.}
\]
Iam si pro $x^\tau$ scribatur $x$, erit
\[
x^m+A'x^{m-1}+B'x^{m-2}+\text{etc.}=(x-a^\tau)(x-b^\tau)(x-c^\tau)\ldots
\]
adeoque
\[
x^m+A'x^{m-1}+B'x^{m-2}+\text{etc.}=0
\]
aequatio quaesita. Quod vero hic $A',B'$ etc. sint non solum rationales sed etiam integri, facile ex theoria aequationis $x^\tau=1$ deducitur.

Quoniam hac operatione in sequentibus saepe utemur, per $(P,\rho^\tau)$ 
\SourceLeaf{234}
indicabimus functionem, qua cifrae aequali posita aequatio proveniens habeat radices, quae sunt potestates $\tau^{\mathrm{tae}}$ radicum aequationis $P=0$.

Si $P\equiv Q$ secundum modulum quemcunque, erit etiam $(P,\rho^\tau)\equiv(Q,\rho^\tau)$ secundum eundem modulum.

\Article{349}
\textsc{Theorema.} \emph{Coëfficiens termini $x^n$ in $(P,\rho^\tau)$ congruus est secundum modulum $\tau$ coëfficiens termini $x^{\tau n}$ in $P^\tau$, siquidem $\tau$ est numerus primus} (quod pro hoc casu est tertia solutio problematis praecedentis).

\emph{Demonstr.} Ex capite sexto sequitur, producti
\[
(x^m+Ax^{m-1}+\text{etc.})(x^m+A\theta x^{m-1}+\text{etc.})\ldots
\]
coëfficientem quemcunque habere hanc formam, postquam pro $\theta^\tau$ substituta est unitas,
\[
E+(1+\theta+\theta\theta+\text{etc.}+\theta^{\tau-1})F
\]
Quodsi iam $\theta$ consideretur tamquam radix prima aequationis $x^\tau=1$, totum productum abiit in $E$; si vero ponatur $\theta=1$, totum productum abiit in $P^\tau=E+\tau F$, quare erit coëfficiens termini $x^{n\tau}$ in $P^\tau$ congruus secundum modulum $\tau$ coëfficienti termini $x^{n\tau}$ in $E$, i. e. coëfficienti termini $x^n$ in $(P,\rho^\tau)$.

\Article{350}
\textsc{Theorema.} \emph{Si $\tau$ est numerus primus, erit}
\[
(P,\rho^\tau)\equiv P\;(\mathit{mod.}\,\tau)
\]

\emph{Demonstr.} Sit coëfficiens termini $x^n$ in $(P,\rho^\tau)=N'$, in $P$ vero eiusdem termini coëfficiens $=N$. Tunc posito
\[
P=x^m+Ax^{m-1}+\text{etc.}+Nx^n+\text{etc.}
\]
erit
\[
P^\tau\equiv x^{m\tau}+A^\tau x^{(m-1)\tau}+\text{etc.}+N^\tau x^{n\tau}+\text{etc.}\;(\mathrm{mod.}\,\tau)
\]
adeoque (§. praec.) $N'\equiv N^\tau\;(\mathrm{mod.}\,\tau)$; quare, quum $N^\tau\equiv N$, erit $N'\equiv N$. Q.\,E.\,D.

Hinc etiam patet, esse $(P,\rho^\alpha)\equiv(P,\rho^{\alpha\tau})$ et $(P,\rho^\tau)\equiv(P,\rho^{\tau\tau})$, unde generaliter
\[
(P,\rho^\alpha)\equiv(P,\rho^{\alpha\tau^k})\;(\mathrm{mod.}\,\tau)
\]

\SourceLeaf{235}
\Article{351}
\textsc{Theorema.} \emph{Datur valor numeri $\nu$ minor quam $p^m$, ita ut functio $x^\nu-1$ per functionem propositam $P$ $m$ dimensionum, cuius pars infima indeterminatam $x$ non involvit, secundum modulum $p$ dividi possit.}

\emph{Dem.} Dividatur per $P$ series functionum $1,x,xx\ldots$ usque ad $x^{p^m-1}$, simulac dimensionem $m$ superant, et quoniam nulla per $P$ sine residuo dividi poterit, omnia residua ad hanc formam redigi poterunt
\[
Ax^{m-1}+Bx^{m-2}+\dotstwo+N
\]
ita ut omnes coëfficientes sint positivi et $<p$. Sed patet, quum nunquam omnes possint esse $=0$, $p^m-1$ tantummodo functiones dari, quarum alicui singulae aequales esse debent, quare quum usque ad potestatem ipsius $x$, cuius exponens est $p^m-1$, $p^m$ residua habeantur, necessario duo ad minimum eadem esse debent. Prodeat igitur idem residuum, si $x^a$ et $x^{a+\nu}$ per $P$ dividantur, ita ut $a+\nu<p^m$. Quare $x^{a+\nu}-x^a$ per $P$ dividi poterit. Hinc quoniam (hyp.) $x$ adeoque etiam $x^a$ functio est ad $P$ prima, etiam $x^\nu-1$ per $P$ dividi poterit Q.\,E.\,D.

\emph{Coroll.} Si $x^\nu-1$ per $P$ dividatur, etiam $x^{k\nu}-1$ per $P$ dividi poterit, denotante $k$ numerum quemcunque integrum.

\Article{352}
\textsc{Theorema.} \emph{Manentibus denominationibus ut in §. praec., si $P$ fuerit functio prima et $x^\nu$ infima potestas, quae unitate mulctata per $P$ dividi possit, erit $\nu$ aut $=p^m-1$ aut pars aliquota huius numeri, excepto unico casu, ubi $P=x$.}

\emph{Dem.} Quoniam $P$ est functio prima $m$ dimensionum, dabuntur $p^m-1$ functiones diversae pauciorum quam $m$ dimensionum (exclusa scilicet ab omnium numero functione $0$), quae omnes ad $P$ erunt primae. Iam quum $x^\nu$ supponatur esse infima potestas, quae per $P$ divisa unitatem relinquit, palam est, si omnes inferiores potestates ab $1,x,\dotstwo$ usque ad $x^{\nu-1}$ per $P$ dividantur, $\nu$ residua diversa prodire, quae per $A$ generaliter designentur. Iam si haec exhauriant omnia quae sunt possibilia, theorema erit demonstratum; sin vero quaedam nondum sint in eorum numero, sit quodcunque eorum $B$; iam perspicuum est, functionem $Bx^\nu$ per $P$ divisam residuum $B$ dare et generaliter esse $Bx^{\nu+k}\equiv Bx^k(\mathrm{mod.}\,P)$; sed omnes functiones ab $B$ usque ad $Bx^{\nu-1}$ diversae inter se et ab residuis $A$
\SourceLeaf{236}
dabunt residua; si scilicet esset $Bx^\lambda\equiv Bx^{\lambda+\delta}(\mathrm{mod.}\,P)$, foret etiam $1\equiv x^\delta(\mathrm{mod.}\,P)$, et $\delta<\nu$ contra hyp.; si vero esset $Bx^\lambda\equiv x^\mu(\mathrm{mod.}\,P)$, foret $B\equiv x^{\mu+\nu-\lambda}(\mathrm{mod.}\,P)$ adeoque $B$ unum ex residuis $A$ contra hyp. Quare patet haberi adhuc $\nu$ nova residua. Simili modo ulterius progredi licebit (omnino ut supra \textsection.\ .) apparebitque numerum omnium residuorum possibilium $p^m-1$ esse aut $=\nu$, aut $=2\nu$, aut $=3\nu$, aut generaliter multiplum numeri $\nu$. Q. E. D.

\Article{353}
Ex theoremate praec. et Coroll. \textsection.\ 351 sequitur, quamvis functionem primam $n$ dimensionum metiri functionem $x^{p^n-1}-1$ secundum modulum $p$. Omnes itaque functiones unius dimensionis excepta unica, quae est $\equiv x$, metientur $x^{p-1}-1$, quod est theorema \textsc{Fermatianum}; omnes autem functiones primae secundi gradus i. e. formae $xx+Ax+B$ metientur functionem $x^{pp-1}-1$ etc. Iam sint numeri $n$ divisores omnes $n,\delta,\delta',\delta''$ etc.\ .\ $1$, patetque, $p^n-1$ etiam per $p^\delta-1$, $p^{\delta'}-1$, $p^{\delta''}-1$ etc. $p-1$ dividi posse, quare functio $x^{p^n-1}-1$ per \emph{omnes functiones primas dimensionum $n,\delta,\delta',\delta''$ etc. usque ad functiones primas unius dimensionis (exclusa functione $x$)} dividi poterit, quare etiam (quum omnes hae functiones sint absolute adeoque etiam inter se primae) per \emph{productum ex omnibus}. Sed idem hoc productum habet $p^n-1$ dimensiones (\textsection.\ 347.) (ob deficientiam unius functionis $x$); quare patet, \emph{hoc productum ipsum ipsi $x^{p^n-1}-1\;(\mathrm{mod.}\,p)$ congruum esse debere.}

\Article{354}
\textsc{Theorema.} \emph{Si functio $x^\nu-1$ per functionem $P$ dividitur, erit}
\[
(P,\rho^{k\nu+t})\equiv(P,\rho^t)
\]
\emph{denotantibus $k,t$ numeros quoscunque integros.}

\emph{Dem.} Sit
\[
P=x^m+Ax^{m-1}+Bx^{m-2}+\text{etc.}
\]
notum est, si
\[
\frac{mx^{m-1}+(m-1)Ax^{m-2}+\text{etc.}}{x^m+Ax^{m-1}+\text{etc.}}
\]
in seriem infinitam formae
\[
m\frac1x+\alpha\frac1{xx}+\beta\frac1{x^3}+\gamma\frac1{x^4}+\text{etc.}
\]
\SourceLeaf{237}
evolvatur, fore $\alpha$ summam radicum aequationis $P=0$, $\beta$ summam quadratorum etc. Unde sine labore deducitur, potestatum $\nu+1$, $\nu+2$ etc.$^{\mathrm{tarum}}$ summam congruam esse summae radicum, quadratorum etc. Hinc vero nisi modulus est aequalis aut inferior numero dimensionum functionis $P$, sequitur esse
\[
(P,\rho^{\nu+1})\equiv P,\quad(P,\rho^{\nu+2})\equiv(P,\rho^2),\quad(P,\rho^{\nu+3})\equiv(P,\rho^3)\text{ etc.}
\]
Istum autem casum infra considerabimus.

\Article{355}
\textsc{Theorema.} \emph{Si in serie}
\[
(P,\rho^0),\quad(P,\rho)\quad(P,\rho^2),\quad(P,\rho^3)\text{ etc.}
\]
\emph{post terminum $\nu^{\mathrm{tum}}$ sequentes primis deinceps sunt congrui, $x^\nu-1$ per $P$ dividi poterit, siquidem $P$ nullum factorem pluries contineat.}

\emph{Dem.} Posito $\frac{\mathrm{d}P}{\mathrm{d}x}=Q$, erit $Q$ functio ad $P$ prima. Sit
\[
\frac QP\equiv\frac Ax+\frac B{xx}+\frac C{x^3}+\text{etc.}
\]
tum post terminum $\frac N{x^\nu}$ sequetur (hyp.)
\[
\frac A{x^{\nu+1}}+\frac B{x^{\nu+2}}+\frac C{x^{\nu+3}}+\text{etc.}
\]
Quare erit
\[
\frac QP\equiv\frac{Ax^{\nu-1}+Bx^{\nu-2}+\text{etc.}}{x^\nu-1}
\]
unde patet, functionem $x^\nu-1$ per $P$ dividi posse. Q. E. D.
